% ===========================================================================
%  第 1 章 引言
% ===========================================================================
\section{引言}

\subsection{背景与动机}

纵列双发矢量推力飞行器（tandem twin-engine vector-thrust aircraft）是一类以
纵列布置的双电机、正交单轴摆座与反向旋转螺旋桨为主要执行机构的垂直起降
（VTOL）概念飞行器。其飞控固件在 2026 年 8 月已完成实机首飞验证并进入调参优化
阶段，姿态控制、执行器分配、遥测链路等子系统均已形成稳定代码基线。在这些子系统
之上，状态估计——尤其是组合导航系统——承担着向控制律提供连续、一致、可用的姿态、
速度与位置信号的职责，是闭环控制、定高/定点模式以及 GNSS 丢失后安全退化共同依赖
的基础设施。

固件中的组合导航系统并非单一滤波器，而是一组按``精度层级''与``任务优先级''
组织的多级状态估计架构：
\begin{itemize}
  \item \textbf{主滤波器}：15 状态误差状态扩展卡尔曼滤波器（Ekf15State），
    状态包括 NED 位置误差、NED 速度误差、体系姿态误差、加速度计零偏误差与陀螺
    零偏误差，以 200\,Hz 的 IMU 捷联惯导解算为时间更新，以 GNSS 位置/速度量测及
    多类辅助量测为量测更新；
  \item \textbf{独立垂直滤波器}：2 状态（高度+速度）卡尔曼滤波器，融合 IMU 垂直
    加速度与气压计/激光测距高度，为垂直控制通道提供与主滤波器相互独立的冗余估计；
  \item \textbf{独立水平滤波器}：北/东两轴各 2 状态的卡尔曼滤波器，融合 IMU 水平
    加速度与光流速度，为无 GNSS 环境下的水平定位提供备份；
  \item \textbf{双矢量航向融合}：利用 GPS 地速矢量与光流机体速度矢量几何解算真实
    航向，作为主滤波器的标量航向量测，抑制陀螺积分航向漂移；
  \item \textbf{辅助逻辑}：统一静止检测器、静止辅助强度分级调度、GNSS 动态权重、
    GNSS 历元时间映射、气压高度基准换算等平台无关纯函数模块。
\end{itemize}

这套架构在工程上体现了``单一高精度主滤波器 + 分层冗余备份 + 防御性门控''的设计
思想：主滤波器保证 GNSS 可用时的最优融合精度；独立滤波器与备用 AHRS 保证 GNSS
丢失时的可用性；NIS 门控、物理残差门限、协方差稳定化与输入合法性防御则保证在
传感器异常、野值、信号失锁等恶劣条件下的数值安全。

本报告的目的，是对这套组合导航系统进行逐层、逐式、逐行的理论与算法分析：
把固件代码中的每一个矩阵块、每一处增益调度、每一组门限常量，还原为严格的数学
推导与明确的物理语义，并基于固件真实参数构建独立数值仿真对其行为进行定量验证。
该工作服务于三个层面的需求：
\begin{enumerate}
  \item \textbf{设计与评审}：为固件参数调优、后续扩展（如 RTK、视觉、气压计零偏
    状态化）提供经得起推敲的数学模型基础；
  \item \textbf{教学与传承}：使 3968 行的滤波器核心代码不再依赖``注释即文档''，
    而具备可与公式一一对照的完整理论映射；
  \item \textbf{验证与审计}：通过独立仿真复现滤波器在典型场景下的行为，检验算法
    语义与实现的一致性，并明确系统的适用边界。
\end{enumerate}

\subsection{分析对象与代码范围}

本报告的分析对象为固件源码目录 \file{TandemVec\_FCS} 下的导航相关文件，
\cref{tab:code-scope} 列出全部被分析文件及其角色。其中 \file{ekf\_15\_state.h}
是滤波器核心（约 3968 行），本报告对其内部数学实现进行逐段还原；\file{navigation\_task.cpp}
是 200\,Hz 导航任务主循环（约 1769 行），本报告对其状态机、辅助调度与输出桥接
进行流程级分析；其余 \file{include/} 下的 \code{ins\_*}.h 系列为平台无关纯函数
模块，可独立编译与宿主测试，本报告对其算法语义逐一给出数学描述。

\begin{table}[H]
\centering
\caption{本报告分析固件代码范围一览}
\label{tab:code-scope}
\small
\begin{tabularx}{\textwidth}{@{}L{5.2cm}L{2.6cm}X@{}}
\toprule
\textbf{文件} & \textbf{规模} & \textbf{角色} \\
\midrule
\file{lib/navigation-main/src/ekf\_15\_state.h} & 3968 行 & 15 状态误差状态 EKF 核心：
  初始化对准、单/双子样时间更新、协方差传播与稳定化、GNSS/速度/重力/姿态/航向/
  静止陀螺/气压高度量测更新、GNSS 延迟回放重传播 \\
\file{src/navigation\_task.cpp} & 1769 行 & 200\,Hz 导航任务：IMU 降采样与双子样
  组装、静止检测调度、GNSS 历元队列消费与质量检查、辅助量测调度、输出桥接、垂直
  速度多源估计、垂直/水平 KF 任务 \\
\file{include/VerticalKF\_2State.h} & 311 行 & 2 状态垂直 KF（高度+速度）\\
\file{include/HorizontalKF.h} & 245 行 & 单轴 2 状态水平 KF（位置+速度，光流速度观测）\\
\file{include/ins\_gnss\_dynamic\_weight.h} & 125 行 & GNSS 动态权重 R（clamp + pDOP 缩放）\\
\file{include/ins\_gnss\_epoch\_timing.h} & 94 行 & GNSS iTOW$\rightarrow$本地时刻映射与量测年龄 \\
\file{include/ins\_static\_detector.h} & 259 行 & 静止检测器（阈值+迟滞+置信度）\\
\file{include/ins\_static\_aid\_profile.h} & 161 行 & 静止辅助强度分级调度 \\
\file{include/ins\_altitude\_reference.h} & 37 行 & 气压高度基准换算 \\
\bottomrule
\end{tabularx}
\end{table}

\subsection{相关工作与技术背景}

组合导航（INS/GNSS 融合）经过数十年发展已形成成熟的理论体系，本报告的分析
建立在以下经典工作的基础之上：

\begin{itemize}
  \item \textbf{捷联惯导编排}：Savage \cite{Savage1998,Savage1998b} 系统地
    建立了捷联惯导算法的统一框架，包括圆锥/划摇补偿的通用形式与误差分析；
    Titterton \& Weston \cite{Titterton2014} 给出工程导向的完整教材性推导；
  \item \textbf{INS/GNSS 融合架构}：Groves \cite{Groves2013} 对松耦合、
    紧耦合与深耦合架构、误差状态滤波、量测建模进行了全面论述；Farrell \&
    Barth \cite{Farrell2000} 侧重 GPS 与惯导融合的实践细节；
  \item \textbf{误差状态四元数滤波}：Sol{\`a} \cite{Sola2017} 给出了
    ES-EKF 中四元数误差状态、指数映射与误差传播的现代系统化表述，与固件
    ``体系右乘误差角''的选取一致；Markley \cite{Markley1999} 分析了不同
    姿态误差表示对滤波性能的影响；
  \item \textbf{姿态滤波与互补滤波}：Madgwick \cite{Madgwick2011} 与 Mahony
    \cite{Mahony2008} 的梯度下降/互补滤波器被固件用作备用 AHRS（\code{backup\_ahrs}）
    与无 GNSS 时的姿态量测源；
  \item \textbf{随机过程与估计理论}：Maybeck \cite{Maybeck1979} 给出随机过程
    模型（高斯-马尔可夫）与卡尔曼滤波的经典论述；Crassidis \& Junkins
    \cite{crassidis2011} 覆盖可观测性（Gramian）与最优估计的理论工具；
  \item \textbf{大地参考}：WGS-84 椭球定义与 Somigliana 重力模型
    \cite{NAVSTAR2000} 是地球模型部分的权威依据。
\end{itemize}

本报告在以上理论框架下，以固件源码为``规范实现''，将每一项通用理论还原为
该特定实现的具体形式（含数值常量、调度策略与门限），形成一份可逐行对照的
工程分析文档。

\subsection{报告组织结构}

本报告沿``从物理到实现、从推导到验证''的递进结构组织，共 11 章：

\begin{itemize}
  \item \textbf{第 2 章（数学基础）}：定义坐标系、四元数与旋转矢量、斜对称算子、
    高斯-马尔可夫随机过程、离散化与线性化等全篇共用的数学工具，并给出误差状态
    卡尔曼滤波的一般框架；
  \item \textbf{第 3 章（地球模型）}：推导 WGS-84 椭球几何（子午圈/卯酉圈曲率
    半径）、Somigliana 正常重力模型、地球自转与传输速率、经纬高（LLA）运动学，
    并逐项对照固件 \file{earth\_model.h} 与 \file{ekf\_15\_state.h} 中的实现；
  \item \textbf{第 4 章（捷联惯导编排）}：推导姿态/速度/位置的离散递推，重点展开
    双子样圆锥补偿与划摇补偿的数学模型、中点重力与科氏力补偿、导航系转动补偿，
    以及固件中单/双子样路径的等价关系；
  \item \textbf{第 5 章（15 状态误差状态 EKF）}：建立 15 维误差状态、连续时间
    系统矩阵 $\m{F}$ 的完整推导（含完整 $\m{F}$ 矩阵扩展项）、协方差传播
    （一阶/二阶 $\bm{\Phi}$ 与 Simpson 离散过程噪声）、Joseph 协方差更新、
    误差状态注入与协方差稳定化；
  \item \textbf{第 6 章（量测模型）}：推导 GNSS 位置/速度量测（含
    LLA$\rightarrow$NED 新息）、速度量测、重力方向量测、姿态量测、标量航向量测、
    静止陀螺量测与气压高度量测的观测矩阵与观测噪声模型，并对每一处与固件实现
    对照；
  \item \textbf{第 7 章（静止辅助与调度）}：推导统一静止检测器的阈值、迟滞状态机
    与置信度评分，静止辅助三档强度调度（ZUPT/Gravity/StaticGyro）的降频策略及其
    可观性依据；
  \item \textbf{第 8 章（鲁棒性机制）}：分析 NIS 门控、物理残差门限、GNSS 失联
    协方差膨胀、低增益重捕获、GNSS 延迟回放重传播、输入合法性防御与协方差数值
    稳定化（Gershgorin/LDLT/特征值投影）的数学模型与设计权衡；
  \item \textbf{第 9 章（任务集成）}：给出 200\,Hz 导航任务的状态机、GNSS 历元
    队列消费、动态权重配置、重锚定逻辑、输出桥接（EKF/备用 AHRS/DETA100 三路
    切换）以及独立垂直/水平滤波器的任务分工；
  \item \textbf{第 10 章（仿真验证）}：基于固件真实参数构建独立数值仿真，设计六个
    场景并给出图表化结果与定量分析；
  \item \textbf{第 11 章（结论）}：总结全篇结论、固件设计特点与后续工作建议。
\end{itemize}

\subsection{主要结论预览}

本报告在分析中形成以下关键结论，正文将给出完整推导与证据：
\begin{enumerate}
  \item 固件采用\textbf{体系右乘误差角} $\delta\vp{\beta}^b$（
    $q = q_{nom}\otimes\Exp(\delta\vp{\beta}^b)$）的姿态误差定义，由此
    $\m{F}_{\Phi V}$ 块带 $\m{C}_n^b$ 旋转项（固件第 2397--2401 行），与
    导航系左乘误差的 KF-GINS 差一个旋转，实为同构；
  \item 协方差传播使用\textbf{双精度 Simpson 离散过程噪声}
    $Q_d = \frac{dt}{6}\left(Q_c + 4\Phi_{\frac12}Q_c\Phi_{\frac12}\T +
    \Phi Q_c\Phi\T\right)$（固件第 2592--2596 行），在 15$\times$15 规模下
    兼顾精度与嵌入式开销；
  \item 静止辅助的量测噪声与频率按置信度\textbf{分三档递减}，其核心依据是
    速度/重力方向量测对陀螺零偏的弱可观性——需要足够长的静止驻留与适当的量测
    频率配合才能收敛；
  \item GNSS 失联后位置协方差按经验增长率膨胀（水平 1.2\,m/s、垂直 1.8\,m/s），
    重捕获时对物理残差合理的量测采用\textbf{膨胀 R 低增益融合}，避免单帧跳变
    （固件第 1620--1644 行）；
  \item GNSS 延迟回放（15\,ms 历史缓冲重传播）将延迟量测严格对齐到其观测历元，
    消除``当前时刻直接融合''引入的时间失配误差。
\end{enumerate}
